\section{Introduction}

This report presents a comprehensive security analysis of a data hosting server that integrates three advanced features: user-controlled security labeling, policy-based data sharing, and untrusted program execution. The analysis employs the \textbf{Volpano type system method} for secure information flow analysis to ensure that the system prevents both confidentiality and integrity violations.

\subsection{Motivation}

Modern data hosting services face growing security challenges:

\begin{itemize}
    \item \textbf{Data sensitivity}: Users store increasingly sensitive data (medical records, financial information, personal photos).
    \item \textbf{Complex sharing requirements}: Data protection often depends on context-specific policies (friends-only photos, private medical records).
    \item \textbf{Untrusted computation}: Users want to run programs on hosted data (filters, anonymization, analysis), but programs may contain bugs or malicious code.
    \item \textbf{Compliance needs}: Organizations (hospitals, financial institutions) require formal security guarantees that policies are enforced.
\end{itemize}

\subsection{Project Scope}

This project designs and verifies a data hosting server that:

\begin{enumerate}
    \item Allows users to assign security labels to data (e.g., \texttt{private}, \texttt{friends}, \texttt{public}).
    \item Enforces label-based sharing policies (e.g., only friends can access friend-labeled data).
    \item Safely executes untrusted programs on user data without leaking information beyond policy.
    \item Provides formal guarantees against information flow violations.
\end{enumerate}

\subsection{Analysis Method: Volpano Type System}

Rather than monitoring program execution at runtime, this project uses the \textbf{Volpano method}---a static type system for secure information flow. Key advantages:

\begin{itemize}
    \item \textbf{Static verification}: Security properties are checked at design/compile time, not runtime.
    \item \textbf{Type-based enforcement}: Security types (Low, High) prevent unauthorized information flows before execution.
    \item \textbf{Sound guarantees}: If type checking succeeds, the program provably preserves security policy.
    \item \textbf{Prevents implicit leaks}: Type rules catch information leaks through control structures, not just direct assignments.
\end{itemize}

\subsection{Report Structure}

\begin{itemize}
    \item \textbf{Section 2}: System participants, roles, and security goals (confidentiality and integrity).
    \item \textbf{Section 3}: Design rationale and Volpano information flow analysis.
    \item \textbf{Section 4}: Server API specification and command definitions.
    \item \textbf{Section 5}: Security analysis of each API command using the Volpano method.
    \item \textbf{Section 6}: Declassification justification for intentional information release.
    \item \textbf{Section 7}: Team contributions and authorship.
    \item \textbf{Section 8}: References and resources.
\end{itemize}

\subsection{Hospital Scenario}

As a motivating example, we consider a hospital using this service:

\begin{itemize}
    \item \textbf{Data}: Medical images (X-rays, CT scans) and patient records.
    \item \textbf{Labels}: Private (only patient), Friends (doctors involved in care), Public (anonymized statistics).
    \item \textbf{Programs}: Image processing (enhancement, analysis), anonymization (removing patient identifiers), statistical analysis.
    \item \textbf{Guarantee}: Even if a program contains malicious code, it cannot leak private medical data to public outputs without explicit declassification.
\end{itemize}

